%%%%%%%%%%%%%%%%%
% CV tailored for Thales cortAIx - Data Scientist / Recherche Opérationnelle
%%%%%%%%%%%%%%%%%

\documentclass[8pt,a4paper,ragged2e,withhyper]{../_shared/altacv}

\geometry{left=1.25cm,right=1.25cm,top=.5cm,bottom=1.cm,columnsep=1.2cm}

\usepackage{paracol}
\usepackage{fontawesome5}
\usepackage{hyperref}

\iftutex 
  \setmainfont{Roboto Slab}
  \setsansfont{Lato}
  \renewcommand{\familydefault}{\sfdefault}
\else
  \usepackage[rm]{roboto}
  \usepackage[defaultsans]{lato}
  \renewcommand{\familydefault}{\sfdefault}
\fi

\definecolor{SlateGrey}{HTML}{2E2E2E}
\definecolor{LightGrey}{HTML}{666666}
\definecolor{DarkPastelRed}{HTML}{8AC0D9}
\definecolor{PastelRed}{HTML}{00B0DA}
\definecolor{GoldenEarth}{HTML}{B4AD0C}

\colorlet{name}{black}
\colorlet{tagline}{PastelRed}
\colorlet{heading}{DarkPastelRed}
\colorlet{headingrule}{GoldenEarth}
\colorlet{subheading}{PastelRed}
\colorlet{accent}{PastelRed}
\colorlet{emphasis}{SlateGrey}
\colorlet{body}{LightGrey}

\definecolor{violet}{HTML}{5A2582}
\definecolor{rouge}{HTML}{F53B3B}
\definecolor{noir}{HTML}{00232C}
\definecolor{bleu}{HTML}{00B0DA}
\definecolor{rose}{HTML}{F831A0}
\definecolor{jaune}{HTML}{FFEC00}
\definecolor{vert}{HTML}{34AD6C}

\renewcommand{\namefont}{\Huge\rmfamily\bfseries}
\renewcommand{\personalinfofont}{\footnotesize}
\renewcommand{\cvsectionfont}{\LARGE\rmfamily\bfseries}
\renewcommand{\cvsubsectionfont}{\large\bfseries}
\renewcommand{\cvItemMarker}{{\small\textbullet}}
\renewcommand{\cvRatingMarker}{\faCircle}

\input{../_shared/pubs-num.tex}
\addbibresource{../_shared/sample.bib}

\begin{document}

\name{BOILEAU Guillaume}
\tagline{Data Scientist / AI Engineer | ML, Signal Processing, MLOps-Aware Workflows \& Critical Systems}

\photoL{2.8cm}{../_shared/moi}

\personalinfo{%
  \email{guillaume.boileaupro@gmail.com}
  \phone{+33 6 59 94 85 84}
  \location{Alpes-Maritimes, FRANCE}
  \linkedin{guillaume-boileau-phd-891b81265}
  \researchgate{Guillaume-Boileau-2}
  \orcid{0000-0002-3576-6968}
}

\makecvheader
\vspace{-0.3cm}

\AtBeginEnvironment{itemize}{\small}
\columnratio{0.64}

\begin{paracol}{2}

\cvsection{Experience}

\cvevent{\textbf{Software Engineer - AI-Aware Systems, Integration \& Production Diagnostics}}{VEV Platform Services France}{2025 -- 2026}{Valbonne, France}

\textbf{Context:} Development and integration of software services for an EV charging CPMS platform connected to operational infrastructure, data flows and hardware systems.

\textbf{Objective:} Support reliable software integration, operational data-flow continuity, system diagnostics and production follow-up for services connected to field equipment.

\begin{itemize}
  \item \textbf{Software development:} Developed and maintained web and backend services using TypeScript, Angular and Node.js in an operational industrial environment.
  \item \textbf{System integration:} Contributed to the integration and maintenance of software services interacting with operational infrastructure and connected equipment.
  \item \textbf{Production diagnostics:} Analysed service behaviour, data exchanges, communication issues and anomalies affecting operational continuity.
  \item \textbf{Data flows:} Implemented, monitored and corrected FTP-based operational data exchange services.
  \item \textbf{Issue analysis:} Investigated incidents involving software behaviour, operational data exchanges and hardware-connected systems.
  \item \textbf{AI systems exposure:} Worked with chatbot and MCP-oriented architectures, with understanding of tool interaction, model integration and application-level AI workflows.
  \item \textbf{Coordination:} Worked with technical stakeholders to clarify issues, follow priorities and support iterative delivery.
  \item \textbf{Tooling:} Used Zenhub, Git-based workflows and collaborative documentation practices for delivery follow-up and technical coordination.
\end{itemize}

\divider

\cvevent{\textbf{Senior R\&D Engineer - Data Science, Signal Processing, ML Validation \& Technical Coordination}}{CNES, Laboratoire Lagrange, Observatoire de la Côte d'Azur}{2023 -- 2025}{Nice, France}

\textbf{Context:} Engineering and data-analysis activities for the LISA ESA/NASA space mission, involving complex software chains, signal analysis, simulation workflows, validation campaigns and international coordination.

\textbf{Objective:} Develop, integrate and validate Python-based data science workflows under strong consistency, traceability, performance and verification constraints.

\begin{itemize}
  \item \textbf{Data science workflows:} Developed Python-based workflows for model integration, comparison, validation and statistical analysis in a mission-level scientific context.
  \item \textbf{Software platform:} Developed CATpyLISA, a Python platform for large-scale model integration, comparison, validation and technical review.
  \textcolor{DarkPastelRed}{\href{https://gitlab.oca.eu/gboileau/lisa_l3_tools}{\bf GitLab: CATpyLISA}}
  \item \textbf{Statistical modelling:} Applied statistical methods, signal-processing techniques and uncertainty analysis to complex low-SNR data environments.
  \item \textbf{Algorithm validation:} Compared model outputs, simulation results and analysis chains to identify deviations, limitations and performance constraints.
  \item \textbf{ML-oriented analysis:} Used Python scientific computing workflows and ML-aware methods to structure model comparison, anomaly detection and performance interpretation.
  \item \textbf{Critical systems context:} Worked in a high-rigor space mission environment involving traceability, technical evidence, review logic and mission-level performance objectives.
  \item \textbf{Technical studies:} Analysed sensitivity to modelling assumptions, numerical parameters, data-processing choices and configuration changes.
  \item \textbf{Validation logic:} Structured comparison campaigns, consistency checks, acceptance logic and validation evidence for technical outputs.
  \item \textbf{Technical reporting:} Produced analysis notes, validation summaries, technical documentation and review material for contributors and project stakeholders.
  \item \textbf{Coordination:} Coordinated contributors, supervised interns and supported technical activities across multiple stakeholders in an international environment.
\end{itemize}

\divider

\cvevent{\textbf{Data Scientist - Statistical Modelling, Noise Diagnostics \& Performance Analysis}}{Universiteit Antwerpen, Belgium}{2021 -- 2023}{Antwerp, Belgium}

\textbf{Context:} Data analysis and performance studies for precision instruments in a high-requirement international environment linked to gravitational-wave detector performance.

\textbf{Objective:} Identify, characterize and mitigate factors affecting system sensitivity using statistical modelling, signal processing and validation-oriented methods.

\begin{itemize}
  \item \textbf{Statistical modelling:} Developed Python-based pipelines for noise source identification, uncertainty estimation, Bayesian inference and statistical diagnostics.
  \item \textbf{Signal processing:} Analysed instrumental and environmental effects affecting measurement quality, sensitivity and stability.
  \item \textbf{Anomaly detection:} Investigated abnormal behaviours, dominant contributors, weak points and possible mitigation directions.
  \item \textbf{Performance analysis:} Identified contributors to performance degradation and supported prioritisation of correction or mitigation actions.
  \item \textbf{Data-driven decisions:} Analysed measurement and simulation outputs to extract meaningful conclusions from complex technical datasets.
  \item \textbf{Reproducibility:} Produced traceable, reproducible and documented analyses for international scientific and engineering stakeholders.
  \item \textbf{Communication:} Communicated complex results through structured reporting, presentations and collaborative review.
\end{itemize}

\divider

\cvevent{\textbf{Junior R\&D Engineer - Signal Analysis, Simulation \& Algorithmic Modelling}}{ARTEMIS, Observatoire de la Côte d'Azur}{2018 -- 2021}{Nice, France}

\textbf{Context:} Data-analysis and simulation activities for gravitational-wave mission preparation, involving weak-signal detection, modelling, simulation and noise characterization.

\textbf{Objective:} Develop robust analysis methods for complex signals and support technical investigations in demanding system environments.

\begin{itemize}
  \item \textbf{Signal analysis:} Developed methods for the separation of overlapping low-SNR signals in complex datasets.
  \item \textbf{Simulation workflows:} Built simulation approaches for complex signal environments, uncertainty studies and large-scale datasets.
  \item \textbf{Algorithm assessment:} Compared simulated outputs, model behaviours and analysis results to assess consistency, robustness and performance limitations.
  \item \textbf{Technical investigations:} Investigated noise sources through cross-sensor correlation analyses in diagnostic contexts.
  \item \textbf{Scientific computing:} Developed strong expertise in Python-based analysis, modelling, documentation and reproducible workflows.
\end{itemize}

\switchcolumn

\cvsection{Profile for Thales cortAIx}

\textbf{Data Scientist / AI Engineer with strong experience in statistical modelling, signal processing, Python-based data science, simulation workflows and algorithm validation for complex and critical systems}. Background developed through major space and gravitational-wave collaborations, complemented by recent industrial software experience involving operational data flows, software integration and production diagnostics.

Experienced in \textbf{Python}, \textbf{statistics}, \textbf{Machine Learning}, \textbf{TensorFlow}, \textbf{PyTorch}, \textbf{Java}, \textbf{GitLab}, \textbf{Docker}, scientific computing, signal processing, anomaly analysis, model comparison, validation workflows and technical reporting. Also familiar with \textbf{chatbot architectures}, \textbf{MCP-based tool interaction}, AI agents, tool-calling logic and model integration into application workflows.

For a \textbf{Data Scientist / Operational Research} role in a critical Defence and Aerospace AI environment, I bring a strong analytical profile: analysing complex data, selecting relevant models, developing and validating algorithms, communicating results clearly and contributing to AI solutions under traceability, security and operational constraints.

\cvsection{Fit for Data Science / AI Delivery}

\begin{itemize}
  \item Strong background in statistics, Bayesian inference, signal processing, uncertainty analysis and complex data interpretation.
  \item Hands-on experience with Python-based data science workflows, scientific computing, model comparison and validation logic.
  \item Experience with TensorFlow, PyTorch, Java, GitLab, Docker and reproducible software environments.
  \item Ability to develop, test and validate algorithms using statistical, mathematical and Machine Learning methods.
  \item Experience analysing weak signals, noisy datasets, anomalies, degradation patterns and performance limitations.
  \item Familiarity with AI system architectures, chatbot systems, MCP, tool-calling logic and model integration into application workflows.
  \item Comfortable communicating complex technical results to mixed technical and non-technical stakeholders.
  \item Strong fit for defence, aerospace and critical-system contexts requiring rigor, autonomy, technical evidence and clear reporting.
\end{itemize}

\cvsection{Team Management}

\begin{itemize}
  \item Supervision of \textbf{3 Master's thesis interns (M2)} during software, modelling and data-analysis activities.
  \item Coordination of technical activities involving \textbf{research engineers and technicians} on a \textbf{data processing and validation pipeline} for the \textbf{LISA space mission}.
  \item Experience organizing tasks, supporting contributors, reviewing outputs and maintaining alignment across collaborative technical developments.
\end{itemize}

\cvsection{Languages}

\cvskill{English}{4}
\divider
\cvskill{Spanish}{2}
\divider

\medskip

\cvsection{Education}

\cvevent{PhD in Physics}{Université Côte d'Azur}{2018 -- 2021}{Nice, France}

\divider

\cvevent{Selective Graduate Program in Fundamental Physics (Magistère)}{Université Paris-Saclay}{2016 -- 2018}{Orsay, France}

\divider

\cvevent{MSc in Astronomy and Astrophysics}{Observatoire de Paris / Université Paris-Saclay}{2017 -- 2018}{Paris, France}

\divider

\cvevent{BSc in Fundamental Physics}{Université Paris-Saclay}{2015 -- 2016}{Orsay, France}

\end{paracol}

\cvsection{Professional Training}

\cvevent{Supervised Machine Learning (Regression \& Classification)}{DeepLearning.AI - Andrew Ng}{2020}{Coursera}

\divider

\cvevent{Recherche reproductible : principes méthodologiques pour une science transparente}{FUN MOOC -- Inria}{2025}{France Université Numérique}

\divider

\cvevent{Continuous Professional Training -- Software, Data, AI and Systems}{OpenClassrooms}{2024 -- 2026}{}

\small{
Python, SQL, Machine Learning, AI, Linux, JavaScript, TypeScript, Angular, Node.js, Django, REST APIs, C/C++, Java, algorithms and software engineering fundamentals.
}

\cvsection{Projects}

\cvevent{CNES / LISA -- Data Science, Signal Processing \& Algorithm Validation}{ESA/NASA Space Mission - Large-scale International Collaboration}{2023 -- 2025}{Nice, France}

Contribution to Python-based data science, signal-processing, model-comparison and validation workflows for the LISA space mission within a large international engineering and scientific collaboration.

\textbf{Data science relevance:} Statistical modelling, uncertainty analysis, weak-signal analysis, simulation workflows, model comparison, anomaly investigation and validation evidence.

\textbf{AI delivery relevance:} Development of reusable Python workflows, GitLab-based software platform, reproducible environments, technical documentation and structured review material.

\textbf{Critical systems relevance:} Mission-level performance analysis, traceability, technical evidence, rigorous validation and work in a high-constraint international aerospace context.

\divider

\cvevent{Einstein Telescope / LIGO--Virgo--KAGRA -- Statistical Diagnostics \& Detector Performance}{International Collaborations}{2021 -- 2023}{Antwerp, Belgium}

Contribution to collaborative developments in data analysis, signal characterization, statistical diagnostics and performance investigations within large-scale international scientific and engineering collaborations.

\textbf{Role:} Development of analysis tools, contribution to technical studies, reporting and collaborative review in high-standard environments.

\textbf{System impact:} Studies on environmental and instrumental noise sources supporting the identification of performance limitations and design constraints for next-generation gravitational-wave observatories.

\cvsection{Strengths}

\vspace{-.3cm}

\cvsubsection{Data science, ML \& AI}

{\small
\cvtag{Data Science}
\cvtag{Machine Learning}
\cvtag{Deep Learning}
\cvtag{Statistical Modelling}
\cvtag{Bayesian Inference}
\cvtag{MCMC}
\cvtag{Regression}
\cvtag{Classification}
\cvtag{Feature Engineering}
\cvtag{Model Selection}
\cvtag{Model Validation}
\cvtag{Algorithm Design}
\cvtag{Operational Research}
\cvtag{Optimization}
\cvtag{Decision Support}
\cvtag{Explainable AI Awareness}
\cvtag{Ethical AI Awareness}
}

\divider\smallskip
\vspace{-.3cm}
\newline
\cvsubsection{Signal, detection \& complex data}

{\small
\cvtag{Signal Processing}
\cvtag{Weak Signal Detection}
\cvtag{Low-SNR Analysis}
\cvtag{Noise Modelling}
\cvtag{Anomaly Detection}
\cvtag{Pattern Recognition}
\cvtag{Data Fusion Awareness}
\cvtag{Simulation Workflows}
\cvtag{Synthetic Data}
\cvtag{Uncertainty Analysis}
\cvtag{Sensitivity Analysis}
\cvtag{Performance Assessment}
\cvtag{Complex Data Analysis}
\cvtag{Scientific Computing}
}

\divider\smallskip
\vspace{-.3cm}

\cvsubsection{AI engineering \& integration}

{\small
\cvtag{AI System Architecture}
\cvtag{Model Integration}
\cvtag{AI Agents}
\cvtag{MCP}
\cvtag{Chatbot Architecture}
\cvtag{LLM Integration}
\cvtag{Tool-Calling}
\cvtag{Prompt Engineering}
\cvtag{Application Workflows}
\cvtag{API Integration}
\cvtag{Data Pipelines}
\cvtag{Reproducible Workflows}
\cvtag{MLOps Awareness}
\cvtag{Deployment Awareness}
\cvtag{User-Facing AI Systems}
}

\divider\smallskip
\vspace{-.3cm}

\cvsubsection{Programming, frameworks \& tools}

{\small
\cvtag{Python}
\cvtag{TensorFlow}
\cvtag{PyTorch}
\cvtag{Java}
\cvtag{Pandas}
\cvtag{NumPy}
\cvtag{Scikit-Learn}
\cvtag{SQL}
\cvtag{Bash}
\cvtag{Linux}
\cvtag{Git}
\cvtag{GitLab}
\cvtag{GitHub}
\cvtag{Docker}
\cvtag{CI/CD}
\cvtag{Jupyter}
\cvtag{REST API}
\cvtag{OpenSearch}
\cvtag{Confluence}
\cvtag{JIRA}
}

\divider\smallskip
\vspace{-.3cm}

\cvsubsection{Critical systems \& defence context}

{\small
\cvtag{Critical Systems}
\cvtag{Defence AI Interest}
\cvtag{Aerospace Context}
\cvtag{Space Systems}
\cvtag{Mission-Level Performance}
\cvtag{Embedded AI Awareness}
\cvtag{Frugal AI Awareness}
\cvtag{Cybersecurity Awareness}
\cvtag{Sensitive Data Awareness}
\cvtag{System Validation}
\cvtag{Technical Evidence}
\cvtag{Traceability}
\cvtag{Risk Identification}
\cvtag{Robustness}
\cvtag{Reliability}
}

\divider\smallskip
\vspace{-.3cm}

\cvsubsection{Collaboration, delivery \& communication}

{\small
\cvtag{Technical Workshops}
\cvtag{Client-Facing Communication}
\cvtag{Scientific Communication}
\cvtag{Technical Reporting}
\cvtag{Knowledge Sharing}
\cvtag{Stakeholder Alignment}
\cvtag{Cross-Functional Coordination}
\cvtag{Agile Environment}
\cvtag{Technical English}
\cvtag{Autonomy}
\cvtag{Initiative}
\cvtag{Problem Solving}
\cvtag{Analytical Thinking}
\cvtag{Pedagogy}
\cvtag{Innovation Mindset}
}

\end{document}
